\section{Introduction}
\label{sec:introduction}

Chemical translation requires preserving distinctions in names, materials, processes, formulas,
and quantities. General MT metrics assess overlap or learned similarity
\cite{papineni2002,post2018,rei2020,zhang2020}; these differ from checking whether a specific technical
term is retained correctly. We investigate how reference-side terminology annotations complement
conventional evaluation, and how providing such terminology affects a controlled translation comparison.

The primary domain is chemistry patent abstracts. JRC-Acquis legal articles offer a distinct,
terminology-sensitive comparison domain, not additional chemical training data. The evaluation unit
is a bounded segment, potentially containing several sentences, rather than a complete patent or
long legal document. The intended scale is a few hundred tokens; actual corpus lengths are reported
instead of assuming a shared 256-token mean.

The contributions are multilingual benchmark construction, provenance-preserving terminology
extraction and curation, and an evaluation protocol separating general quality from term retention.
These preserve the chemistry-aware evaluation objective. The infrastructure is implemented, but
the full translation comparison remains planned; no model ranking or superiority claim is made.